\documentclass[11pt]{article}
\usepackage[a4paper,margin=1in]{geometry}
\usepackage{amsmath,amssymb,amsthm,booktabs,enumitem,microtype}
\usepackage[numbers,sort&compress]{natbib}
\usepackage{cleveref}

\newtheorem{theorem}{Theorem}[section]
\newtheorem{corollary}[theorem]{Corollary}
\newtheorem{proposition}[theorem]{Proposition}
\newtheorem{lemma}[theorem]{Lemma}
\theoremstyle{remark}
\newtheorem{remark}[theorem]{Remark}
\newcommand{\Tr}{\operatorname{Tr}}
\newcommand{\cB}{\mathcal B}
\newcommand{\cC}{\mathcal C}
\newcommand{\cG}{\mathcal G}
\newcommand{\cP}{\mathcal P}
\newcommand{\cS}{\mathcal S}
\newcommand{\one}{\mathbf 1}

\title{Projector-Mass First-Alias Threshold\\
and an Isospectral Cell Obstruction}
\author{Bin Wang}
\date{August 2026}

\begin{document}
\maketitle

\begin{abstract}
The projector-density gauge of RH-335 turns the global noisy parity scalar
into signed localized cells.  We determine the exact first-alias scale of one
such cell.  On every bounded-phase natural clock, its parity contribution
divided by $H_k=kR^{-2k}$ equals
$2C_*\lambda^{\eta_\sigma}(\beta R)^{2k}\pi_\sigma(J)(1+o(1))$.
Thus the critical projector-mass exponent is
$\kappa_{\rm proj}=\log(\beta R)/\log\lambda$, distinct from the earlier
Duhamel stability exponent.  A fixed finite partition necessarily has a maximum
normalized parity cell diverging to $+\infty$, but the maximizing cell may
move with noise and is not identified with a physical boundary/sibling
aggregate; raw and alias terms may still cancel it.  We complement the scale
theorem with an exact three-state family.  A nontrivial interval of smooth
similarities preserves strict positivity, row-stochasticity, the spectrum,
and every power trace, while the parity projector masses and corrected
singleton cells drift by explicit nonzero zero-sum vectors.  This is a
nonphysical finite algebraic obstruction to recovering cells from global
spectral data.  It does not prove physical signed Duhamel cancellation or a
physical normalized obstruction, and it yields no determinant or
Riemann-hypothesis conclusion.
\end{abstract}

\section{Inherited projector gauge and first-alias clock}

Let $K_\sigma$ be the noisy folded operator and let
$\lambda_-(\sigma)=-(1-\delta_\sigma)$ be its real simple parity
eigenvalue, where the archived law is
\begin{equation}\label{eq:delta-law}
 \delta_\sigma=C_*\sqrt\sigma+o(\sqrt\sigma),
 \qquad C_*>0.
\end{equation}
Write $E_{-,\sigma}$ for the rank-one Riesz projector itself.  RH-335 proves
that
\begin{equation}\label{eq:pi}
 \pi_\sigma(J)=\Tr(M_JE_{-,\sigma})
\end{equation}
is a real finite signed measure with total mass one
\citep{WangProjectorLedger2026}.  Allocating the deterministic parity scalar
with this noisy density is a frozen gauge, not canonical physical
localization or deterministic/noisy projector transport.

Retain the constants and clock from RH-326
\citep{WangFirstAlias2026}:
\begin{equation}\label{eq:constants}
 r_H=\frac{17}{20},\qquad R=\frac75,\qquad
 \beta=\frac1{r_H\sqrt\lambda},
\end{equation}
and
\begin{equation}\label{eq:phase}
 \eta_\sigma
 =k-\frac{\log(1/\sigma)}{2\log\lambda}.
\end{equation}
Throughout the moving statements, $k=k_\sigma\in\mathbb N$, $k\ge2$,
$\eta_\sigma$ remains bounded, and $k\to\infty$ as $\sigma\to0$.

For a measurable cell $J$, possibly depending on $\sigma$, define
\begin{equation}\label{eq:G-H}
 \cG_{\sigma,k}(J)
 =r_H^{-2k}\{1-\lambda_-(\sigma)^{2k}\}\pi_\sigma(J),
 \qquad H_k=kR^{-2k}.
\end{equation}
This is only the localized parity-gauge constituent of the RH-335 corrected
cell.  It omits the localized raw defect and the globally subtracted alias
packet.

\section{Moving projector-mass threshold}

\begin{theorem}[Exact first-alias projector-mass scale]
\label{thm:moving-scale}
On every bounded-phase sequence,
\begin{equation}\label{eq:G-over-H}
 \boxed{
 \frac{\cG_{\sigma,k}(J)}{H_k}
 =2C_*\lambda^{\eta_\sigma}(\beta R)^{2k}
 \pi_\sigma(J)\{1+o(1)\}.}
\end{equation}
The scalar $o(1)$ is independent of the chosen cell.
\end{theorem}

\begin{proof}
Because $k\delta_\sigma\to0$, the uniform parity expansion of RH-326 gives
\begin{equation}\label{eq:parity-expansion}
 1-(1-\delta_\sigma)^{2k}
 =2kC_*\sqrt\sigma\{1+o(1)\}.
\end{equation}
Even order removes the sign of $\lambda_-$, so substitution in
\eqref{eq:G-H} and division by $H_k$ yield
\[
 2C_*\sqrt\sigma\left(\frac{R}{r_H}\right)^{2k}
 \pi_\sigma(J)\{1+o(1)\}.
\]
Equation \eqref{eq:phase} gives
$\sqrt\sigma\lambda^k=\lambda^{\eta_\sigma}$, while
$(\beta R)^{2k}=(R/r_H)^{2k}\lambda^{-k}$.  Their product is the required
scalar in \eqref{eq:G-over-H}.
\end{proof}

Define
\begin{equation}\label{eq:kappa}
 \boxed{
 \kappa_{\rm proj}
 =\frac{\log(\beta R)}{\log\lambda}
 =\frac{\log(28/17)}{\log\lambda}-\frac12.}
\end{equation}
The rational bracket $3/2<u_c<25/16$ from the physical observation source
implies
$\lambda=2u_c(u_c-1)<225/128<(28/17)^2$
\citep{WangObservation2026}.  Hence $\beta R>1$ and
$\kappa_{\rm proj}>0$.  Ordinary substitution of the archived decimal for
$\lambda$ gives
\begin{equation}\label{eq:kappa-decimal}
 \kappa_{\rm proj}=0.463406944517002\ldots .
\end{equation}
This decimal is diagnostic, not an interval certificate.  It is not the
RH-325 sufficient Duhamel weight ceiling
$\gamma_*=0.3503698834605293\ldots$; in fact their separation is exact.
The RH-334 algebraic identities imply
$\lambda^3+4\lambda^2-16=0$ \citep{WangObservation2026}.  This polynomial is
strictly increasing on the positive half-line, and its value at $17/10$ is
$473/1000>0$, so $\lambda<17/10$.  Moreover,
\[
 \left(\frac{196}{85}\right)^2-\left(\frac{17}{10}\right)^3
 =\frac{116783}{289000}>0.
\]
Since $\lambda>1$ and
\[
 \kappa_{\rm proj}-\gamma_*
 =\frac{\log\!\left((R^2/r_H)/\lambda^{3/2}\right)}{\log\lambda},
 \qquad \frac{R^2}{r_H}=\frac{196}{85},
\]
it follows that $\kappa_{\rm proj}>\gamma_*$ exactly.  The two exponents
also govern different objects.

\begin{corollary}[Negligibility, exact phase conversion, and critical limit]
\label{cor:threshold}
On every bounded-phase sequence,
\begin{equation}\label{eq:iff}
 \cG_{\sigma,k}(J)=o(H_k)
 \quad\Longleftrightarrow\quad
 \pi_\sigma(J)=o((\beta R)^{-2k}).
\end{equation}
Moreover,
\begin{equation}\label{eq:phase-conversion}
 \boxed{
 (\beta R)^{-2k}
 =\sigma^{\kappa_{\rm proj}}
  (\beta R)^{-2\eta_\sigma}.}
\end{equation}
If
\begin{equation}\label{eq:critical-assumption}
 (\beta R)^{2k}\pi_\sigma(J)\longrightarrow p,
 \qquad \eta_\sigma\longrightarrow\eta,
\end{equation}
then
\begin{equation}\label{eq:critical-limit}
 \frac{\cG_{\sigma,k}(J)}{H_k}
 \longrightarrow 2C_*\lambda^\eta p.
\end{equation}
\end{corollary}

\begin{proof}
Bounded phase makes $2C_*\lambda^{\eta_\sigma}\{1+o(1)\}$ bounded above
and away from zero.  This proves \eqref{eq:iff} from
\cref{thm:moving-scale}.  From \eqref{eq:phase},
$k=\log(1/\sigma)/(2\log\lambda)+\eta_\sigma$; inserting this into the
left side of \eqref{eq:phase-conversion} proves the exact identity.
Finally apply \eqref{eq:critical-assumption} to \eqref{eq:G-over-H}.
\end{proof}

\section{A forced large parity cell, but not a physical obstruction}

\begin{theorem}[Fixed-partition maximum]
\label{thm:partition-max}
Let $\cP=\{J_1,\ldots,J_N\}$ be one fixed finite measurable partition.
Then
\begin{equation}\label{eq:max-pi}
 \max_{1\le i\le N}\pi_\sigma(J_i)\ge\frac1N.
\end{equation}
Consequently,
\begin{equation}\label{eq:max-G}
 \max_{1\le i\le N}
 \frac{\cG_{\sigma,k}(J_i)}{H_k}\longrightarrow+\infty.
\end{equation}
The maximizing index may depend on $\sigma$.  Along every sequence
$\sigma_\ell\to0$, some fixed index recurs as a maximizer on a subsequence,
and its normalized parity cell diverges on that subsequence.
\end{theorem}

\begin{proof}
RH-335 gives
$\sum_{i=1}^N\pi_\sigma(J_i)=1$, so the maximum is at least the arithmetic
mean $1/N$, even though individual masses may be negative.  The scalar in
\eqref{eq:G-over-H} is positive for small noise.  Since bounded phase makes
$\lambda^{\eta_\sigma}$ uniformly bounded below and $\beta R>1$,
\eqref{eq:max-pi} proves \eqref{eq:max-G}.  The subsequence statement is the
pigeonhole principle on the finite index set.
\end{proof}

\begin{remark}[Cancellation firewall]
The theorem concerns only the parity-gauge constituent.  It does not
identify a maximizing cell with the physical $\cB+\cS$ aggregate, and it
does not prove a physical nonzero normalized obstruction.  The localized
raw defect can cancel the parity term inside a corrected cell, other cells
can cancel in the partition sum, and the first-alias coefficient further
subtracts $\mathcal A_{k,2k}$.  The maximizing index may also move with
noise.  Thus physical signed $\Delta_B+\Delta_S$ cancellation remains
\texttt{NOT\_TESTABLE}.
\end{remark}

\section{A strictly positive isospectral Markov family}

We now give a separate finite algebraic witness.  Reuse the exact RH-335
matrix and parity projector
\begin{align}
 K&=\begin{pmatrix}
 3/17&7/51&35/51\\
 4/85&83/255&32/51\\
 58/85&1/51&76/255
 \end{pmatrix},\label{eq:K}\\
 E_-&=\begin{pmatrix}
 10/17&-5/51&-25/51\\
 8/17&-4/51&-20/51\\
 -10/17&5/51&25/51
 \end{pmatrix}.\label{eq:E}
\end{align}
For $t\ne1$, put
\begin{equation}\label{eq:S-family}
 S_t=\begin{pmatrix}1-t&t&0\\0&1&0\\0&0&1\end{pmatrix},
 \qquad K_t=S_t^{-1}KS_t.
\end{equation}
Define $q(t)=35-38t-12t^2$.  Exact multiplication gives the audited formula
\begin{equation}\label{eq:Kt-formula}
 K_t=\begin{pmatrix}
 \dfrac{15-4t}{85}&\dfrac{q(t)}{255(1-t)}
   &\dfrac{35-32t}{51(1-t)}\\[0.8em]
 \dfrac{4(1-t)}{85}&\dfrac{83+12t}{255}&\dfrac{32}{51}\\[0.8em]
 \dfrac{58(1-t)}{85}&\dfrac{5+174t}{255}&\dfrac{76}{255}
 \end{pmatrix}.
\end{equation}

\begin{theorem}[Positive row-stochastic isospectral interval]
\label{thm:positive-family}
For every
\begin{equation}\label{eq:sufficient-interval}
 \boxed{-\frac5{174}<t<\frac12,}
\end{equation}
the matrix $K_t$ is strictly positive and row-stochastic.  For every such
$t$,
\begin{equation}\label{eq:spectrum}
 \operatorname{spec}(K_t)=\{1,-2/5,1/5\},
\end{equation}
and for every integer $m\ge1$,
\begin{equation}\label{eq:all-traces}
 \boxed{
 \Tr K_t^m=1+(-2/5)^m+(1/5)^m.}
\end{equation}
\end{theorem}

\begin{proof}
On \eqref{eq:sufficient-interval}, every denominator in
\eqref{eq:Kt-formula} is positive.  The nonconstant numerators
\[
 15-4t,\quad q(t),\quad35-32t,\quad83+12t,\quad5+174t
\]
are also positive; for $q$, it suffices that $t<1/2$ because its positive
root exceeds $1/2$.  More directly, $q'(t)=-38-24t<0$ on the stated
interval and $q(1/2)=13$, so $q(t)>13$.  Thus every entry is positive.  Since
$S_t\one=\one$, also $S_t^{-1}\one=\one$, and
$K_t\one=S_t^{-1}K\one=\one$.

Similarity preserves the spectrum and gives
$K_t^m=S_t^{-1}K^mS_t$.  RH-335 supplies the three distinct eigenvalues of
$K$, so trace invariance proves \eqref{eq:all-traces}.
\end{proof}

\begin{remark}[Sufficient versus maximal]
The interval \eqref{eq:sufficient-interval} is convenient, not maximal.
Intersecting all strict entry inequalities in the connected component
containing zero gives
\begin{equation}\label{eq:maximal-interval}
 \left(-\frac5{174},\frac{-19+\sqrt{781}}{12}\right).
\end{equation}
The upper endpoint is the positive root of $q(t)$.  No maximality claim is
made beyond this explicit finite family.
\end{remark}

\section{Moving projector masses and corrected singleton cells}

Let
\begin{equation}\label{eq:Et}
 E_-(t)=S_t^{-1}E_-S_t.
\end{equation}
It is the Riesz projector of $K_t$ at $-2/5$ by similarity.

\begin{proposition}[Exact projector-mass drift]
\label{prop:pi-drift}
For the three singleton windows,
\begin{equation}\label{eq:pi-t}
 \boxed{
 \pi(t)=\operatorname{diag}E_-(t)
 =\left(\frac{10-8t}{17},\frac{-4+24t}{51},\frac{25}{51}\right).}
\end{equation}
Therefore
\begin{equation}\label{eq:pi-drift}
 \pi(t)-\pi(0)=\frac{8t}{17}(-1,1,0),
 \qquad \sum_i\pi_i(t)=1.
\end{equation}
\end{proposition}

\begin{proof}
Insert \eqref{eq:E} and \eqref{eq:S-family} into \eqref{eq:Et} and read the
diagonal.  Subtraction gives \eqref{eq:pi-drift}; its entries sum to zero.
\end{proof}

Use now the RH-335 fixed-order fixture $n=2$, $r_H=17/20$, parity
eigenvalue $-2/5$, and zero deterministic singleton slots.  Define
\begin{equation}\label{eq:Ct-def}
 \cC_i(t)=\left(\frac{20}{17}\right)^2
 \left\{(K_t^2)_{ii}+\frac{21}{25}\pi_i(t)\right\}.
\end{equation}

\begin{theorem}[Isospectral corrected-cell obstruction]
\label{thm:C-drift}
The corrected singleton cells are
\begin{equation}\label{eq:C-t}
 \boxed{
 \cC(t)=\left(
 \frac{6800-5760t}{4913},
 \frac{400+5760t}{4913},
 \frac{6672}{4913}\right).}
\end{equation}
For every admissible $t$,
\begin{align}
 \sum_i\cC_i(t)&=\frac{48}{17},\label{eq:C-total}\\
 \cC(t)-\cC(0)&=\frac{5760t}{4913}(-1,1,0).
 \label{eq:C-drift}
\end{align}
At $t=1/100$, the drift is exactly
\begin{equation}\label{eq:one-percent-drift}
 \left(-\frac{288}{24565},\frac{288}{24565},0\right).
\end{equation}
\end{theorem}

\begin{proof}
Similarity applied to $K^2$ gives
\[
 \operatorname{diag}K_t^2
 =\left(\frac{215-192t}{425},
 \frac{53+192t}{425},\frac{242}{425}\right).
\]
Combine this with \eqref{eq:pi-t} in \eqref{eq:Ct-def} and simplify to
\eqref{eq:C-t}.  Summation and subtraction prove
\eqref{eq:C-total}--\eqref{eq:C-drift}; setting $t=1/100$ gives
\eqref{eq:one-percent-drift}.
\end{proof}

The full spectrum and all power traces are constant, yet the corrected
singleton cells move.  This proves nonidentification of those cells from
global spectral data in a strictly positive row-stochastic family.
RH-210 already proved the general logical phenomenon that projectors may
move under a fixed divisor \citep{WangDivisorPivot2026}.  The narrow addition
here is the simultaneous positivity, Markov normalization, all-power trace
lock, and explicit corrected-cell ledger.

The family remains nonphysical finite algebra.  In particular, fixed
$n=2$ is not a $k=1$ first-alias counterloop: the archived counterloop
definition requires $k\ge2$.

\section{Reproduction and claim boundary}

The artifact evaluates every matrix identity with exact rational arithmetic.
It checks the displayed formulas, the exact rational exponent-separation
certificate, strict positivity at rational sample
points inside the proved interval, the endpoint factorization, two-sided
projector intertwining, twelve power traces, projector masses, corrected
cells, invariant sums, and the exact one-percent drift.  Floating-point rows
only reproduce the symbolic phase conversion and the diagnostic exponent;
they are not interval certificates.

RH-336 proves the moving projector-mass threshold and a nonphysical
isospectral corrected-cell obstruction.  It does not identify a physical
$\Delta_B+\Delta_S$ cancellation, prove a physical nonzero normalized
obstruction, exclude raw/local/alias signed cancellation, close the far or
off-alias terms, transport the noisy head, or glue a determinant.  No result
constructs a Hilbert--Polya operator, identifies Riemann zeros, proves a von
Mangoldt trace formula or completed-zeta divisor equality, or proves the
Riemann hypothesis.  Gates A--E remain false/open.

\bibliographystyle{plainnat}
\bibliography{references}

\end{document}
